% ── §3.3: Αγωγός Προ-επεξεργασίας (Preprocessing Pipeline) ──

\section{Αγωγός Προ-επεξεργασίας (Preprocessing Pipeline)}
\label{sec:preprocessing}

Η προ-επεξεργασία μετατρέπει τα ακατέργαστα σήματα αισθητήρων σε
τυποποιημένη μορφή κατάλληλη για εισαγωγή στον 1D-CNN ταξινομητή.
Ο αγωγός αποτελείται από τρία διαδοχικά στάδια: αφαίρεση θορύβου,
κανονικοποίηση και τεμαχισμό με συρόμενο παράθυρο. Εφαρμόζεται
ανεξάρτητα σε κάθε client (subject) — κομβική επιλογή για τη
διατήρηση της ατομικής κατανομής στο FL σενάριο.

% ─── 3.3.1 Denoising ────────────────────────────────────────
\subsection{Αφαίρεση Θορύβου}
\label{subsec:denoising}

Τα σήματα επιταχυνσιόμετρου και γυροσκοπίου περιέχουν δύο
συνιστώσες \cite{Anguita2013}:

\begin{itemize}
    \item \textbf{Συνιστώσα βαρύτητας (DC):} Χαμηλής συχνότητας
    σήμα ($< 0{,}3\,\text{Hz}$) που αντικατοπτρίζει την κλίση
    της συσκευής ως προς τη βαρύτητα. Σταθερό κατά τη διάρκεια
    κάθε δραστηριότητας.

    \item \textbf{Συνιστώσα κινητικής επιτάχυνσης (body):}
    Υψηλής συχνότητας σήμα ($> 0{,}3\,\text{Hz}$) που απεικονίζει
    τη δυναμική κίνηση του σώματος. Αυτό είναι το πληροφοριακά
    χρήσιμο σήμα για HAR.
\end{itemize}

Ο διαχωρισμός υλοποιείται με \textbf{φίλτρο Butterworth 3ης
τάξης} με συχνότητα αποκοπής $f_c = 0{,}3\,\text{Hz}$:

\begin{equation}
    \label{eq:butterworth}
    |H(j\omega)|^2 = \frac{1}{1 + \left(\dfrac{\omega}{\omega_c}\right)^{2n}},
    \quad n = 3
\end{equation}

\noindent Το low-pass φίλτρο εξάγει τη συνιστώσα βαρύτητας, ενώ
η αφαίρεσή της από το ολικό σήμα παράγει το σήμα body acceleration
— τα τρία επιπλέον κανάλια ($b_x, b_y, b_z$) που ανεβάζουν το
σύνολο των καναλιών από 6 σε 9 \cite{Anguita2013}.

% ─── 3.3.2 Normalization ────────────────────────────────────
\subsection{Κανονικοποίηση}
\label{subsec:normalization}

Η κανονικοποίηση διασφαλίζει ότι όλα τα κανάλια συνεισφέρουν
ισότιμα στην εκπαίδευση, ανεξάρτητα από τη φυσική μονάδα μέτρησής
τους (m/s² για επιτάχυνση, rad/s για γωνιακή ταχύτητα). Εφαρμόζεται
\textbf{Z-score κανονικοποίηση} ανά κανάλι, υπολογισμένη επί του
συνόλου εκπαίδευσης κάθε client:

\begin{equation}
    \label{eq:zscore}
    \hat{x}_{c,t} = \frac{x_{c,t} - \mu_c}{\sigma_c}
\end{equation}

\noindent όπου $\mu_c$ και $\sigma_c$ ο μέσος όρος και η τυπική
απόκλιση του καναλιού $c$ υπολογισμένα αποκλειστικά στα δεδομένα
εκπαίδευσης του αντίστοιχου client. Η εφαρμογή \textit{ανά client}
— και όχι καθολικά — είναι κρίσιμη: αποφεύγει τη διαρροή
πληροφορίας από το σύνολο ελέγχου (data leakage) και τηρεί την
αρχή του FL ότι δεδομένα δεν κοινοποιούνται μεταξύ clients.

Στο UCI HAR dataset, τα δεδομένα παρέχονται ήδη κανονικοποιημένα
στο $[-1, 1]$, οπότε το βήμα αυτό εφαρμόζεται κυρίως στο AAL
Lab dataset \cite{AAL_UCI2016}.

% ─── 3.3.3 Sliding Window ───────────────────────────────────
\subsection{Τεμαχισμός με Συρόμενο Παράθυρο}
\label{subsec:sliding_window}

Η ακατέργαστη χρονοσειρά κάθε client τεμαχίζεται σε παράθυρα
σταθερού μήκους με τη μέθοδο \textit{sliding window}, παράγοντας
τα δείγματα εισόδου για τον ταξινομητή. Χρησιμοποιούνται οι
παράμετροι του UCI HAR \cite{Anguita2013}:

\begin{itemize}
    \item \textbf{Μήκος παραθύρου:} $W = 128$ δείγματα
    ($= 2{,}56\,\text{s}$ στα 50\,Hz) — επαρκές για να περιέχει
    2–5 πλήρεις κύκλους βηματισμού.

    \item \textbf{Βήμα ολίσθησης:} $S = 64$ δείγματα, αντίστοιχο
    επικάλυψης $\rho = 50\%$ — διπλασιάζει τον αριθμό παραθύρων
    και αυξάνει τη διαφορετικότητα των training samples.
\end{itemize}

Για μια χρονοσειρά $T$ συνολικών δειγμάτων, ο αριθμός παραθύρων
είναι $\lfloor (T - W)/S \rfloor + 1$. Κάθε παράθυρο
$\mathbf{X}^{(k)} \in \mathbb{R}^{9 \times 128}$ αποτελεί ένα
ανεξάρτητο δείγμα εισόδου με ετικέτα $y^{(k)} \in \{1,\ldots,6\}$,
ορισμένη ως η δραστηριότητα που κατέχει την πλειοψηφία των δειγμάτων
του παραθύρου. Ο συνολικός αγωγός προ-επεξεργασίας απεικονίζεται
στο Σχήμα~\ref{fig:preprocessing_pipeline}.

\begin{figure}[H]
    \centering
    \includegraphics[width=\textwidth]{eda_preprocessing_pipeline.pdf}
    \caption{Αγωγός προ-επεξεργασίας από ακατέργαστο σήμα αισθητήρα
    σε τανυστή εισόδου CNN. (1)~Ακατέργαστο σήμα (κίνηση + βαρύτητα).
    (2)~Φίλτρο Butterworth 3ης τάξης ($f_c = 0{,}3$\,Hz): εξαγωγή
    βαρύτητας (πορτοκαλί) και σώματος (κυανό).
    (3)~Z-score κανονικοποίηση ανά client.
    (4)~Sliding window ($W=128$, 50\% επικάλυψη).
    (5)~Τανυστής εξόδου $\mathbf{X} \in \mathbb{R}^{9 \times 128}$.
    Ο αγωγός εφαρμόζεται τοπικά σε κάθε FL client.}
    \label{fig:preprocessing_pipeline}
\end{figure}

Ο αγωγός εφαρμόζεται \textbf{τοπικά σε κάθε client}, σύμφωνα
με τις αρχές του FL: τα στατιστικά κανονικοποίησης ($\mu_c$,
$\sigma_c$) δεν κοινοποιούνται στον server, διατηρώντας την
ιδιωτικότητα των δεδομένων. Η έξοδος κάθε client είναι ένα
τοπικό σύνολο παραθύρων $\{\mathbf{X}^{(k)}, y^{(k)}\}_{k=1}^{K_i}$,
όπου $K_i$ ο αριθμός παραθύρων του client $i$.

Αξίζει να σημειωθεί ότι ο αγωγός αυτός δεν εφαρμόζει μείωση διαστατικότητας (dimensionality reduction) ή επιλογή χαρακτηριστικών (feature selection) — τεχνικές που έχουν εξεταστεί εκτενώς στη βιβλιογραφία~\cite{Zebari2022}. Η επιλογή αυτή γίνεται σκόπιμα: ο 1D-CNN εκτελεί αυτόματη εξαγωγή ιεραρχικών χαρακτηριστικών από το ακατέργαστο σήμα, χωρίς να απαιτείται προηγούμενη χειροκίνητη επιλογή. Η χρήση του συνόλου των 9 καναλιών (3 επιταχύνσεις, 3 γωνιακές ταχύτητες, 3 body accelerations) χωρίς περαιτέρω μείωση επιτρέπει στο δίκτυο να ανακαλύψει από μόνο του ποιοι συνδυασμοί καναλιών είναι πληροφοριακοί για κάθε κλάση δραστηριότητας — προσέγγιση που αποδείχθηκε ανώτερη των hand-crafted features στα πειράματα αξιολόγησης~\cite{Chen2021, Hammerla2016}.
